\section{矩阵元}

由两个胶子场组成的算符（四个指标均不缩并）的核子自旋平均矩阵元由
\begin{align}
{ M}_{\mu \alpha; \lambda \beta } (z,p) \equiv \langle p | \,
G_{\mu \alpha} (z) \, [z,0]\, G_{\lambda \beta} (0) | p \rangle \ ,
\end{align}
给出，其中 $[z,0]$ 是胶子（伴随）表示中的标准直线规范链接
\begin{align}
[x,y]~\equiv~{\rm Pexp}\Big\{ig\!\int_0^1\! dt~(x-y)^\mu \tilde{A}_\mu(tx+(1-t)y)\Big\}
\ .
\label{straightE}
\end{align}
对不变振幅作分解的张量结构可由两个可用四矢量 $p_\alpha$、$z_\alpha$
以及度规张量 $g_{\alpha \beta}$ 构成。考虑到 $G_{\rho \sigma}$ 对其指标的反对称性，我们有
\begin{align}
& { M}_{\mu \alpha; \lambda \beta} (z,p)= \nonumber \\ &
\left( g_{\mu\lambda} p_\alpha p_\beta - g_{\mu\beta} p_\alpha p_\lambda - g_{\alpha\lambda} p_\mu p_\beta + g_{\alpha\beta} p_\mu p_\lambda \right) \mathcal{M}_{pp} \nonumber \\
+& \left( g_{\mu\lambda} z_\alpha z_\beta - g_{\mu\beta} z_\alpha z_\lambda - g_{\alpha\lambda} z_\mu z_\beta + g_{\alpha\beta} z_\mu z_\lambda \right) \mathcal{M}_{zz} \nonumber \\
+ & \left( g_{\mu\lambda} z_\alpha p_\beta - g_{\mu\beta} z_\alpha p_\lambda - g_{\alpha\lambda} z_\mu p_\beta + g_{\alpha\beta} z_\mu p_\lambda \right) \mathcal{M}_{zp} \nonumber \\
+ & \left( g_{\mu\lambda} p_\alpha z_\beta - g_{\mu\beta} p_\alpha z_\lambda - g_{\alpha\lambda} p_\mu z_\beta + g_{\alpha\beta} p_\mu z_\lambda \right) \mathcal{M}_{pz} \nonumber \\
+& \left( p_\mu z_\alpha - p_\alpha z_\mu\right) \left( p_\lambda z_\beta - p_\beta z_\lambda\right) \mathcal{M}_{ppzz}
\nn + & \left(g_{\mu\lambda} g_{\alpha\beta} -g_{\mu\beta} g_{\alpha\lambda} \right)\mathcal{M}_{gg} \ ,
\label{Manb}
\end{align}
其中振幅 ${\cal M}$ 是不变间隔 $z^2$ 与 Ioffe 时间 \cite{Ioffe:1969kf}
$(pz)\equiv - \nu$ 的函数（引入负号是为了后面的方便）。

由于矩阵元应当对场的交换对称（相当于 $\{\mu \alpha\} \leftrightarrow \{\lambda \beta\}$ 且 $z \to -z$），
函数 $\mathcal{M}_{pp}$、$\mathcal{M}_{zz}$、$\mathcal{M}_{gg}$、$\mathcal{M}_{ppzz}$ 与
$\mathcal{M}_{pz} - \mathcal{M}_{zp}$ 是 $\nu$ 的偶函数，而 $\mathcal{M}_{pz} + \mathcal{M}_{zp}$ 关于 $\nu$
为奇。

通常的光锥胶子分布由
$g^{\alpha \beta} { M}_{+ \alpha; \beta +} (z,p)$ 得到，其中 $z$ 取光锥``负''方向，
$z=z_-$。我们有
\begin{align}
g^{\alpha \beta} { M}_{+ \alpha; \beta +} (z_-,p) =-2 p_+ ^2 \mathcal{M}_{pp}(\nu ,0) \ ,
\label{lcpp}
\end{align}
即 PDF 由 ${\cal M}_{pp}$ 结构决定：
\begin{align}
- \mathcal{M}_{pp} (\nu,0) =\frac12
\int_{-1}^1 \dd x \, e^{-i x \nu} x f_g (x) \, \ .
\label{glPDF}
\end{align}
因此，我们应选取其参数化中含有 ${\cal M}_{pp}$
的指标组合 $\{\mu \alpha; \lambda \beta \}$ 的算符。

注意，在此胶子 PDF 定义中自然的量是被胶子携带的动量密度 \mbox{$G(x) \equiv x f_g (x)$}，
而不是其数目密度 $f_g (x)$。
在定域 $z_-=0$（即 $\nu=0$）极限下，该积分给出胶子携带的强子 plus 动量的份额。
在没有胶子--夸克跃迁时，此份额守恒，
这就对 Altarelli-Parisi 核 \cite{Altarelli:1977zs} 的 $gg$ 分量施加了限制：
作用于 $G(x)$ 时它应具有加 prescription（plus-prescription）性质。

由于 $G_{\rho \sigma}$ 对其指标的反对称性，
$\alpha=+$ 与 $\beta= +$ 的值被排除在式 (\ref{lcpp}) 的求和之外。
此外，由于 $g_{- -} =0$，组合
$g^{\alpha \beta} { M}_{+ \alpha; \beta +} (z,p)$ 只包含对横向指标 $i,j =1,2$ 的求和，
即化为 $g^{ij} { M}_{+ i; j +} (z,p) \equiv { M}_{+ i; + i} (z,p)$
（这里改用对 $i$ 的欧氏求和），对此我们有
\begin{align}
{ M}_{+ i; +i}
= { M}_{0 i; 0i} + { M}_{3 i; 3i} + ({ M}_{0 i; 3i} +
{ M}_{3 i; 0i}) \ .
\label{ii}
\end{align}
在定域 $z_3=0$ 极限下，这三个组合分别正比于 ${\bf E}^2_\perp$、${\bf B}^2_\perp$
以及坡印廷矢量的第三分量 $({\bf E \times B})_3$。

这些组合（对 $i$ 求和）在 ${\cal M}$ 结构基底上的分解为
\begin{align}
{ M}_{0 i; i 0 } = & 2 p_0^2 \mathcal{M}_{pp} +2 \mathcal{M}_{gg} \ , \\
{ M}_{3 i; i 3 } = &2 p_3^2 \mathcal{M}_{pp}+2 z_3^2 \mathcal{M}_{zz} \nonumber \\ &
+2 z_3 p_3 \left( \mathcal{M}_{zp} + \mathcal{M}_{pz} \right) - 2\mathcal{M}_{gg} \ ,
\\
{ M}_{0 i; i 3 } = & 2 p_0 \left( p_3 \mathcal{M}_{pp} + z_3 \mathcal{M}_{pz} \right) \ ,
\\
{ M}_{3 i; i 0 } = & 2 p_0 \left( p_3 \mathcal{M}_{pp} + z_3 \mathcal{M}_{zp} \right) \ .
\end{align}

它们都含有定义胶子分布的 $\mathcal{M}_{pp}$ 函数，但运动学因子各不相同。
遗憾的是，没有一个是纯的 $\mathcal{M}_{pp}$：它们都含污染项。
而且，${ M}_{3 i; i 3}$ 矩阵元（最初 \cite{Ji:2013dva}
提议用于从格点提取胶子 PDF）
含有三个污染项，其余的只有一个附加项。
特别地，矩阵元 ${ M}_{0 i; i 0}$ 以 $\mathcal{M}_{gg}$ 为污染项。
容易看出
\begin{align}
{ M}_{j i ; ij} \equiv \bra{p} G_{j i} (z) G_{ij}(0) \ket{p} &= -2 \mathcal{M}_{gg} \ ,
\end{align}
其中假定同时对 $i$ 与 $j$ 求和。因此，组合
\begin{align}
{ M}_{0 i; i 0} + { M}_{j i ; i j } = & 2 p_0^2 \mathcal{M}_{pp} \
\label{00m}
\end{align}
可用于提取扭度 2 函数 $\mathcal{M}_{pp}$。

组合不同类型的矩阵元时，应注意：
在光锥之外，这些矩阵元具有与规范链接的存在相关的额外紫外发散。
由于紫外发散的定域性质，
每个矩阵元无论其指标组合 $\{\mu \alpha; \lambda \beta\}$ 如何，
对这些发散都是乘法可重正的 \cite{Li:2018tpe}。
然而，选取不同的 $\{\mu \alpha; \nu \beta\}$ 组合，一般会得到不同的反常量纲。

因此，这些矩阵元的哪些线性组合是乘法可重正的并非 {\it 先验} 显然。
文献 \cite{Zhang:2018diq} 中证明，式 (\ref{ii}) 所示的组合，即
${ M}_{0 i; i 0}$、${ M}_{3 i; i 3}$、\mbox{${ M}_{0 i; i3} +
{ M}_{3 i; i 0}$}（以及 \mbox{${ M}_{0 i; i 3} -
{ M}_{3 i; i 0}$}），在对横向指标 $i$ 求和后，
各自在单圈水平上都是乘法可重正的。

此外，组合 $G_{ij}G_{ij}$（对横向 $i,j$ 求和）
等于 $2G_{12} G_{12}$，
其矩阵元是乘法可重正的。
如我们将看到的，它与 ${M}_{0 i; i 0}$ 具有相同的单圈 UV 反常量纲，
故式 (\ref{00m}) 的组合在单圈水平上是乘法可重正的。
它在更高阶是否仍然如此，值得进一步研究。

含有对 $\alpha$ 与 $\beta$ 协变求和的组合 $g^{\alpha \beta} { M}_{3 \alpha; 3 \beta}$
也被发现是乘法可重正的。
它由
\begin{align}
& g^{\alpha \beta} { M}_{3 \alpha; 3 \beta } =
\left(2p_3^2 - m^2 \right) \mathcal{M}_{pp} + 3z_3^2 \mathcal{M}_{zz}
\nonumber \\ & + 3 p_3 z_3 \left( \mathcal{M}_{zp}+ \mathcal{M}_{pz} \right) + p_0^2 z_3^2 \mathcal{M}_{ppzz}
-3\mathcal{M}_{gg} \ ,
\end{align}
给出，且污染项最多（四个）。

同样涉及协变求和的函数 $g^{\alpha \beta} { M}_{0 \alpha; 0 \beta}$
曾用于首次尝试 \cite{Fan:2018dxu}
从格点提取胶子 PDF。但如文献
\cite{Zhang:2018diq} 所指出，它不是乘法可重正的。

在任何耦合常数无量纲的理论中，
矩阵元 ${ M}(z,p)$ 都含有对应于微扰演化（或称 DGLAP 演化，即 Dokshitzer-Gribov-Lipatov-Altarelli-Parisi
\cite{Gribov:1972ri,Altarelli:1977zs,Dokshitzer:1977sg}）的 \mbox{$\sim \ln (-z^2)$} 项。
人们或许想知道哪些组合在单圈具有对角的 DGLAP 演化。

为回答这些问题，我们计算了
单圈
胶子交换对原双定域算符的修改。
